\documentclass[aspectratio=169,11pt]{beamer}
\usepackage{iftex}
\ifXeTeX
  \usepackage{fontspec}
  \setsansfont{texgyreheros}[Extension=.otf, UprightFont=*-regular, BoldFont=*-bold,
    ItalicFont=*-italic, BoldItalicFont=*-bolditalic]
\else
  \usepackage[T1]{fontenc}
  \usepackage[utf8]{inputenc}
\fi
\usepackage[turkish,shorthands=off]{babel}
\usepackage{booktabs}
\usepackage{tikz}
\usetikzlibrary{positioning,arrows.meta,fit,backgrounds}

% AUTO-GENERATED by scripts/make_report_assets.py -- do not edit
\newcommand{\GolubN}{72}
\newcommand{\GolubP}{7{,}129}
\newcommand{\GolubMapped}{6{,}434}
\newcommand{\GolubTrainALL}{27}
\newcommand{\GolubTrainAML}{11}
\newcommand{\GolubTestALL}{20}
\newcommand{\GolubTestAML}{14}
\newcommand{\GolubValues}{513{,}288}
\newcommand{\GolubNonFinite}{0}
\newcommand{\GolubNegPct}{28.7}
\newcommand{\GolubBelowPct}{32.7}
\newcommand{\GolubAbovePct}{0.83}
\newcommand{\GolubMin}{-28{,}400}
\newcommand{\GolubMax}{71{,}369}
\newcommand{\GolubAffx}{58}
\newcommand{\GolubAllNeg}{233}
\newcommand{\GolubCallP}{28.2}
\newcommand{\GolubCallA}{70.2}
\newcommand{\GolubCallM}{1.6}
\newcommand{\GolubMedP}{662}
\newcommand{\GolubMedA}{33}
\newcommand{\GolubNeverPresent}{1{,}971}
\newcommand{\GolubPresMin}{19.1}
\newcommand{\GolubPresMax}{40.5}
\newcommand{\GolubFilt}{3{,}859}
\newcommand{\GolubSkewRaw}{7.08}
\newcommand{\GolubSkewLog}{0.55}
\newcommand{\GolubRhoLin}{0.95}
\newcommand{\GolubRhoLog}{0.18}
\newcommand{\GolubPCOne}{10.8}
\newcommand{\GolubPCTwo}{8.8}
\newcommand{\GolubCumTen}{49.6}
\newcommand{\GolubNEighty}{37}
\newcommand{\GolubPCOneAUC}{0.984}
\newcommand{\GolubSilPCA}{0.17}
\newcommand{\GolubSilTSNE}{0.40}
\newcommand{\GolubSig}{945}
\newcommand{\GolubSigOne}{567}
\newcommand{\GolubSigMWU}{1{,}014}
\newcommand{\GolubUpAML}{462}
\newcommand{\GolubUpALL}{483}
\newcommand{\GolubRhoSNR}{0.999}
\newcommand{\GolubARI}{0.94}
\newcommand{\GolubTop}{MGST1, CST3, TCF3, ZYX, CD33, APLP2, CD79B, POU2AF1, DNTT, CCND3}
\newcommand{\GolubPresTrain}{29.0}
\newcommand{\GolubPresTest}{27.3}
\newcommand{\GolubPresP}{0.071}
\newcommand{\GolubQP}{0.016}
\newcommand{\GolubQTrain}{482}
\newcommand{\GolubQTest}{421}
\newcommand{\GolubKSAll}{17.7}
\newcommand{\GolubKSAml}{8.0}
\newcommand{\GolubKSqAll}{28}
\newcommand{\GolubKSqAml}{0}
\newcommand{\AndN}{121}
\newcommand{\AndP}{21{,}370}
\newcommand{\AndMapped}{15{,}936}
\newcommand{\AndMissingPct}{7.19}
\newcommand{\AndInf}{0}
\newcommand{\AndGenesMissing}{12{,}084}
\newcommand{\AndGenesMissingTwenty}{2{,}841}
\newcommand{\AndMaxSampleMissing}{24.3}
\newcommand{\AndMin}{-10.64}
\newcommand{\AndMax}{9.56}
\newcommand{\AndMedian}{-0.04}
\newcommand{\AndNegPct}{52.7}
\newcommand{\AndOutlierPct}{0.81}
\newcommand{\AndPCOne}{17.5}
\newcommand{\AndPCTwo}{8.6}
\newcommand{\AndCumTen}{51.8}
\newcommand{\AndPCOneAUC}{0.999}
\newcommand{\AndBestPC}{1}
\newcommand{\AndBestPCAUC}{0.999}
\newcommand{\AndSilPCA}{0.28}
\newcommand{\AndSilTSNE}{0.41}
\newcommand{\AndSig}{9{,}157}
\newcommand{\AndSigPct}{42.8}
\newcommand{\AndUpAML}{4{,}580}
\newcommand{\AndUpALL}{4{,}577}
\newcommand{\AndARI}{0.41}
\newcommand{\AndTop}{PILRB, INOC1, NPY, NPR1, PCBD1, C17orf58, SPG20, LOC201229, CTSC}
\newcommand{\AndBALL}{87}
\newcommand{\AndTALL}{11}
\newcommand{\AndALL}{98}
\newcommand{\AndAML}{23}
\newcommand{\AndExcluded}{6}
\newcommand{\CommonAll}{\textbf{??}}
\newcommand{\CommonGA}{\textbf{??}}
\newcommand{\CommonGM}{\textbf{??}}
\newcommand{\EtaRawClass}{\textbf{??}}
\newcommand{\EtaRawDataset}{\textbf{??}}
\newcommand{\EtaStdClass}{\textbf{??}}
\newcommand{\EtaStdDataset}{\textbf{??}}
\newcommand{\GenesAnd}{\textbf{??}}
\newcommand{\GenesGolub}{\textbf{??}}
\newcommand{\GenesMile}{\textbf{??}}
\newcommand{\MedAnd}{\textbf{??}}
\newcommand{\MedGolub}{\textbf{??}}
\newcommand{\MedMile}{\textbf{??}}
\newcommand{\MileALL}{\textbf{??}}
\newcommand{\MileAML}{\textbf{??}}
\newcommand{\MileARI}{\textbf{??}}
\newcommand{\MileBestPCAUC}{\textbf{??}}
\newcommand{\MileMax}{\textbf{??}}
\newcommand{\MileMin}{\textbf{??}}
\newcommand{\MileMissingPct}{\textbf{??}}
\newcommand{\MileN}{\textbf{??}}
\newcommand{\MileNSubtypes}{\textbf{??}}
\newcommand{\MileP}{\textbf{??}}
\newcommand{\MileSigPct}{\textbf{??}}
\newcommand{\OvGA}{\textbf{??}}
\newcommand{\OvGAp}{\textbf{??}}
\newcommand{\OvGM}{\textbf{??}}
\newcommand{\OvGMdir}{\textbf{??}}
\newcommand{\OvGMp}{\textbf{??}}
\newcommand{\OvN}{\textbf{??}}
\newcommand{\RhoGA}{\textbf{??}}
\newcommand{\RhoGM}{\textbf{??}}
\newcommand{\RhoMA}{\textbf{??}}
\newcommand{\SharedTop}{\textbf{??}}
\newcommand{\SharedTopN}{\textbf{??}}

\graphicspath{{../../results/figures/}}
\newcommand{\fig}[2][]{\IfFileExists{../../results/figures/#2}{\includegraphics[#1]{#2}}{\fbox{eksik: \texttt{\detokenize{#2}}}}}

\definecolor{ink}{HTML}{0B0B0B}
\definecolor{accent}{HTML}{2A78D6}
\definecolor{aml}{HTML}{EB6834}
\setbeamercolor{frametitle}{fg=accent}
\setbeamercolor{title}{fg=accent}
\setbeamercolor{structure}{fg=accent}
\setbeamercolor{normal text}{fg=ink}
\setbeamertemplate{navigation symbols}{}
\setbeamertemplate{footline}{\hfill\scriptsize\color{gray}\insertframenumber/\inserttotalframenumber\hspace*{4mm}\vspace*{2mm}}
\setbeamerfont{frametitle}{series=\bfseries}

\title{ALL ve AML Lösemi Alt Tip Sınıflandırmasında Sızıntısız Değerlendirme ve Platformlar Arası Aktarım}
\subtitle{Ara Rapor: Literatür, Veri Setleri ve Keşifsel Veri Analizi}
\author{Enes AY \quad 503020260092}
\institute{Eskişehir Osmangazi Üniversitesi, Fen Bilimleri Enstitüsü, Yapay Zeka Programı\\Makine Öğrenmesi ve Anomali Tespiti}
\date{Kasım 2026}

\begin{document}

\begin{frame}
\titlepage
\end{frame}

\begin{frame}{Problem ve motivasyon}
\begin{columns}[T]
\column{0.55\textwidth}
\begin{itemize}
\item ALL (lenfoid) ve AML (miyeloid): ikisi de akut \textbf{lösemi alt tipi}; tedavi protokolleri farklı.
\item Modüldeki ``non-cancer'' etiketi biyolojik olarak yanlış $\Rightarrow$ \emph{ALL vs.\ AML alt tip sınıflandırması}.
\item Mikrodizi: binlerce genin mRNA düzeyini tek deneyde ölçer.
\item Golub: 72 hasta, \GolubP{} prob $\Rightarrow$ $p \gg n$.
\end{itemize}
\column{0.42\textwidth}
\textbf{Araştırma soruları (AS)}
\begin{enumerate}\small
\item Nested CV'de 3 yöntem farklı mı?
\item 38/34 CV ile tutarlı mı? Sızıntı ne kadar şişirir?
\item Sinyal kohortlar arası taşınıyor mu?
\item (Bonus) DL / TabPFN daha iyi mi?
\end{enumerate}
\end{columns}
\end{frame}

\begin{frame}{Literatür taraması}
\begin{columns}[T]
\column{0.48\textwidth}
\begin{itemize}\small
\item Crossref, Europe PMC, PubMed; 6 Kasım'da hedefli ek tarama
\item 47 kaynak tam metin, 4 özet, 1 künye (yapay zekâ destekli okuma)
\item \textbf{52 kaynak}, 31'i 2021+; DOI'ler Crossref ile doğrulandı
\item Sızıntı: Kapoor 2023, Rosenblatt 2024
\item Dış doğrulama: Zheng 2026, Riley 2024; lösemi: Mulet 2025
\end{itemize}
\column{0.5\textwidth}
\textbf{Boşluk}
\begin{itemize}\small
\item Yeni 4 çalışmada \%97--99,9 doğruluk, ama protokoller karşılaştırılamaz
\item Bazı çalışmalarda gen seçimi CV \emph{dışında} $\Rightarrow$ seçim yanlılığı
\item Yeni 4 çalışmada güven aralığı yok; okunan Golub çalışmalarında aktarım yok
\item AUC--eşik ayrışması sepsiste gösterildi; ALL/AML'de bulamadık
\end{itemize}
\end{columns}
\end{frame}

\begin{frame}{Veri setleri}
\centering\small
\begin{tabular}{@{}lccc@{}}
\toprule
 & Golub & MILE (GSE13159) & Andersson (GSE7186) \\
\midrule
Platform & Hu6800 & U133 Plus 2.0 & 27K cDNA (2 renk) \\
Ölçü & MAS4 avg.\ diff. & ölçekli log [0,1] & $\log_2$ oran \\
$n$ (ALL/AML) & 72 (47/25) & \MileN{} (\MileALL/\MileAML) & \AndN{} (\AndALL/\AndAML) \\
$p$ & \GolubP{} & \MileP{} & \AndP{} \\
Eksik değer & 0 & \%\MileMissingPct & \%\AndMissingPct \\
Rol & ana & benzer 1 & benzer 2 (B4) \\
\bottomrule
\end{tabular}

\vspace{2mm}
\fig[width=0.95\textwidth]{fig_class_distribution.pdf}
\end{frame}

\begin{frame}{Veri kalitesi ve dönüşüm (Golub)}
\centering
\fig[width=0.76\textwidth]{fig_golub_quality.pdf}\\[1mm]
\fig[width=0.76\textwidth]{fig_golub_log_distributions.pdf}
\begin{itemize}\small
\item NaN/INF yok; değerlerin \textbf{\%\GolubNegPct}'si negatif (PM $-$ MM), Absent çağrılarıyla (\%\GolubCallA) örtüşüyor.
\item $[20, 16000]$ kırpma + $\log_2$: çarpıklık \GolubSkewRaw{} $\to$ \GolubSkewLog; ortalama--SS korelasyonu \GolubRhoLin{} $\to$ \GolubRhoLog.
\end{itemize}
\end{frame}

\begin{frame}{Sinyal: boyut indirgeme ve ayırt edici genler}
\centering
\fig[width=0.72\textwidth]{fig_golub_embedding.pdf}\\[1mm]
\fig[width=0.72\textwidth,height=0.3\textheight,keepaspectratio]{fig_golub_heatmap.pdf}
\begin{itemize}\small
\item PC1 varyansın \%\GolubPCOne'ini açıklıyor, ama tek başına AUC \GolubPCOneAUC{} $\Rightarrow$ güven aralığı ve test şart.
\item $q<0{,}05$: \GolubSig{} gen. AML: CST3, CD33, MPO; ALL: CD19, CD79B, TCF3.
\item \textcolor{aml}{\textbf{Uyarı:}} sıralama tüm veriyle yapıldı, yalnızca betimsel; modelde kullanılırsa sızıntı.
\end{itemize}
\end{frame}

\begin{frame}{Batch etkisi ve kohortlar arası farklar}
\centering
\fig[width=0.76\textwidth]{fig_golub_batch.pdf}\\[1mm]
\fig[width=0.76\textwidth]{fig_cross_dataset.pdf}
\begin{itemize}\small
\item Sınıf içi KS: ALL'de genlerin \%\GolubKSAll'si $p<0{,}05$ (beklenen \%5) $\Rightarrow$ 34 örneklik test skoru sapabilir.
\item Biyoloji taşınıyor ($\rho=\RhoGM$, ilk \OvN{} genin \OvGM'ü ortak); ölçek taşınmıyor (PC1'de kohort $\eta^2=\%\EtaRawDataset$).
\end{itemize}
\end{frame}

\begin{frame}{Fizibilite pilotu: planlanan ölçümler mümkün mü?}
\begin{columns}[T]
\column{0.40\textwidth}
\small
\textbf{Sızıntı} (38 Golub eğitim örneği, DD)
\begin{itemize}
\item Gerçek etiket: içeride \PilotInReal, dışarıda \PilotOutReal
\item Rastgele etiket: içeride \PilotInPerm, dışarıda \textbf{\PilotOutPerm}
\end{itemize}
\vspace{1mm}
\textbf{Sabit aktarım}
\begin{itemize}
\item \PilotRawSingle/6 yönde AUC \PilotRawAucMin--\PilotRawAucMax, DD 0,50
\item Sepsiste de bildirildi (Zheng 2026)
\item Sıra dönüşümü burada kurtarmıyor: özgün soru
\end{itemize}
\vspace{1mm}
\footnotesize Sabit ayar, ön bulgu; kilitli setlere dokunulmadı.
\column{0.58\textwidth}
\scriptsize
{\setlength{\tabcolsep}{2.5pt}\begin{tabular}{@{}lccccc@{}}
\toprule
Eğitim $\to$ test & AUC ham & DD ham & AUC $z$ & DD $z$ & DD sıra\\ \midrule
Golub $\to$ MILE & 0{,}97 & 0{,}50 & 0{,}97 & 0{,}89 & 0{,}69\\
Golub $\to$ Andersson & 0{,}97 & 0{,}50 & 0{,}97 & 0{,}94 & 0{,}70\\
MILE $\to$ Golub & 0{,}87 & 0{,}84 & 0{,}93 & 0{,}93 & 0{,}64\\
MILE $\to$ Andersson & 1{,}00 & 0{,}92 & 0{,}99 & 0{,}93 & 0{,}52\\
Andersson $\to$ Golub & 0{,}98 & 0{,}50 & 0{,}97 & 0{,}59 & 0{,}50\\
Andersson $\to$ MILE & 0{,}99 & 0{,}50 & 0{,}99 & 0{,}74 & 0{,}50\\
\bottomrule
\end{tabular}
}

\vspace{2mm}
\small\textbf{Kilitli dış testler:} Golub 34 örnek; MILE aşama 2 (693 akut lösemi, özel çip); TARGET (RNA-seq)
\end{columns}
\end{frame}

\begin{frame}{Bulgulardan final planına}
\begin{columns}[T]
\column{0.48\textwidth}
\small
\textbf{Bulgular $\Rightarrow$ kararlar}
\begin{itemize}
\item $p\gg n$ $\Rightarrow$ öznitelik seçimi \textbf{her eğitim katının içinde}
\item Güçlü sinyal $\Rightarrow$ güven aralığı, eşleştirilmiş testler
\item Dengesizlik $\Rightarrow$ tabakalı katlar, dengeli doğruluk
\item Eksik değer $\Rightarrow$ kat içi imputasyon
\item Platform farkı $\Rightarrow$ etiketsiz düzeltmeler ve ortak \CommonAll{} gen
\end{itemize}
\column{0.5\textwidth}
\small
\textbf{Yöntemler:} elastic-net LR, SVM, random forest (+ karar ağacı)\\[1mm]
\textbf{Protokol:} 38/34, test \emph{bir kez}; 72 örnekte 5$\times$5 iç içe CV, 5 tohum\\[1mm]
\textbf{İstatistik:} düzeltilmiş $t$ + Holm, McNemar\\[1mm]
\textbf{Bonuslar:} B1 MLP/otokodlayıcı; B2 TabPFN, GenePT; B3 İngilizce makale (aday IEEE EMBC 2027, hocanın yazılı onayıyla); B4 Andersson\\[1mm]
\textbf{Tekrarlanabilirlik:} \texttt{python run\_all.py}, tohum 42
\end{columns}
\end{frame}

\begin{frame}
\centering\Large Teşekkürler\\[4mm]
\normalsize Kod: \texttt{github.com/mrayenesay/leukemia-all-aml-ml}\\[2mm]
\small GenAI beyanı: Claude (Opus 5.5) kullanıldı; ayrıntılar Ek A'da.
\end{frame}

\end{document}
